% Cover letter for the IP&M submission of the Sieve manuscript.
% All factual statements are drawn from the manuscript (paper-journal,
% 2026-09-24 version) and its auto-generated numbers.tex macros.
\documentclass[12pt]{article}
\usepackage[T1]{fontenc}
\usepackage[margin=1in]{geometry}
\usepackage{microtype}
\usepackage{parskip}
\pagestyle{empty}
\setlength{\parskip}{0.55em}

\newcommand{\PaperTitle}{Sieve: Training-Free, Model-Free, Question-Conditioned
Context Compression for Efficient Long-Context LLM Inference: An
Accuracy--Compression--Cost Frontier Study of Model-Free Information Selection}

\begin{document}

\begin{flushleft}
September 28, 2026

\bigskip
The Editors\\
Information Processing \& Management
\end{flushleft}

\bigskip
\noindent\textbf{Re: Submission of a Full Length Article --- ``\PaperTitle{}''}

\bigskip
Dear Editors,

Please consider the manuscript referenced above for publication in
\emph{Information Processing \& Management} as a \textbf{Full Length Article}.
I am the sole author of the manuscript and serve as its corresponding author.

The manuscript studies a question that sits between two well-mapped regions
of the prompt-compression landscape. Neural compressors such as LLMLingua
and its successors recover most task quality at high ratios, but they
require a hosted scorer, GPU capacity for at least one forward pass over
the full context, and a tokenizer matched to the scorer. Truncation is
free but question-agnostic, and it interacts badly with the positional
biases of long-context models. Between the two lies a region that has not
been mapped systematically: compressors that consult the question, cost
milliseconds of arithmetic, and require no model at all. The manuscript
measures how much quality that region gives up, at what compression, and
at what computational cost.

The vehicle of the study is \emph{Sieve}, a training-free, model-free
sentence selector built from three corpus-local statistics: IDF-weighted
coverage of query content terms, a U-shaped positional prior, and a
redundancy penalty applied greedily under an explicit word budget, with
selected sentences emitted in their original order. Its preprocessing is a
linear-time pass over the context, its budgeted selection is
$O(k\,n\,\bar{s})$, and on the evaluation container the full selector runs
in tens of milliseconds on one CPU core (53.7\,ms at a 6{,}000-word
context), with no model parameters, no training, and no dependency beyond
NumPy.

The evaluation is organized as a three-way accuracy--compression--cost
frontier study over three evidence-annotated benchmarks --- QASPER,
HotpotQA (distractor setting), and 2WikiMultihopQA --- totaling 1{,}600
records under a reader-free evidence-recall protocol, plus end-task
question answering with a production reader at the $4\times$ operating
point, with sample sizes stated and claims scoped accordingly. The main
findings, all reported with paired bootstrap intervals, are: (i) Sieve
retains more gold evidence than every question-agnostic baseline on all
three datasets, with margins over head truncation of 12.8, 44.8, and 22.8
recall points at $4\times$; (ii) a BM25 sentence-ranking baseline, which
the manuscript adds as the strongest question-conditioned model-free
comparator, exceeds the full selector in raw recall on all three
datasets --- a result the manuscript reports as measured and explains
through the selector's own ablations as the positional prior trading
recall for placement where readers attend reliably; (iii) a GPT-2
perplexity gate, measured on the same hardware as every other method,
reaches 28.7\% evidence recall at $4\times$ on QASPER at 255--320$\times$
Sieve's measured cost; and (iv) evidence recall correlates only weakly
with end-task answer F1 (pooled Pearson 0.30), so the two protocols are
reported separately throughout. A formal analysis states the cost of each
pipeline stage and the selection behavior under an explicit
lexical-coverage condition, and a robustness study covers evidence
placement, hyperparameter sensitivity, and cross-dataset transfer without
retuning.

I believe the manuscript fits the scope of \emph{Information Processing \&
Management}. It is a quantitative study of question-conditioned information
selection: it builds directly on classical information-retrieval machinery
(IDF weighting, BM25, maximal marginal relevance), evaluates selection
quality with recall protocols over annotated evidence, and connects
selection decisions to the processing and cost of managing long contexts
in information systems. The frontier analysis reports quality,
compression, and compressor cost for every method, giving practitioners a
measured basis for deciding when statistical selection suffices and when
a hosted scorer earns its infrastructure.

The manuscript substantially extends an earlier conference-format study of
the same method, as described in the Introduction, with a three-dataset
evaluation, a new BM25 baseline, a same-hardware measured comparison
against a neural perplexity gate, formal analysis, and a recall--accuracy
correlation study.

I confirm that: the manuscript is original and has not been published
previously in any form; it is not under consideration by any other
journal; I have no known competing financial interests or personal
relationships that could have appeared to influence the work reported; and
the work received no specific funding. The benchmarks used are public
(QASPER, HotpotQA, 2WikiMultihopQA), and the released repository
reproduces every table and figure from raw checkpoints.

Thank you for considering the manuscript. I look forward to the reviewers'
comments.

\bigskip
Sincerely,

\bigskip
Chang Tan\\
Undergraduate Student, School of Economics and Management\\
Dalian University of Technology\\
No.\,2 Linggong Road, Ganjingzi District, Dalian, Liaoning 116024, China\\
E-mail: 2352935097@qq.com

\end{document}
